% Automatisch aus der Word-/LaTeX-Konvertierung {\"u}bernommen.
\chapter{Methodik und Versuchsaufbau}

In diesem Kapitel wird gezeigt, wie die in Kapitel 1 formulierte Forschungsfrage empirisch untersucht wird, bevor in den nachfolgenden Kapiteln die Implementierung und die Ergebnisse folgen. Zuerst wird das grundlegende Forschungsdesign vorgestellt und in die unabh{\"a}ngige, die abh{\"a}ngigen und die kontrollierten Variablen des Versuchsplans unterteilt (4.1). Anschlie{\ss}end werden die verwendete Zielplattform (4.2) und die Softwareumgebung (4.3) beschrieben. Zudem wird die unabh{\"a}ngige variable, die drei Konfigurationen (4.4) und die f{\"u}nf Laufzeit- und Heap-Benchmarks und den dabei bestimmten statischen RAM-/ROM-Fu{\ss}abdruck, also die eigentlichen Benchmark-Metriken (4.5) erl{\"a}utert. Zuletzt wird eine vollst{\"a}ndige {\"U}bersicht der kontrollierten Variablen (4.6) erstellt, welche {\"u}ber alle Konfigurationsstufen und Metriken hinweg konstant bleibt. Dadurch werden die Reproduzierbarkeit und die Aussagekraft des Konfigurationsvergleichs sichergestellt.

\section{Forschungsdesign}
Die Arbeit verfolgt ein quantitatives und experimentelles Forschungsdesign mit kontrollierter Variation der Zephyr-Konfiguration. Dabei werden wiederholte Messungen auf derselben Hardwareplattform durchgef{\"u}hrt. Dieselbe physische Plattform, ein nrf52840-DK, durchl{\"a}uft nacheinander drei Firmware-Builds, die sich jeweils in der zu untersuchenden Kernel-/ Stack-Konfiguration unterscheiden. Alle {\"u}brigen relevanten Rahmenbedingungen, darunter Hardware, CPU-Taktfrequenz, Zephyr-Version, Compiler und Messverfahren, werden dabei konstant gehalten. Dadurch werden hardwarebedingte Unterschiede zwischen den Messungen minimiert. Zudem k{\"o}nnen Unterschiede in den Messergebnissen auf die Konfigurations{\"a}nderungen zur{\"u}ckgef{\"u}hrt werden, soweit andere Einflussfaktoren kontrolliert werden. Quantitativ bedeutet im Kontext dieser Untersuchung, dass diese anhand von messbaren numerischen Gr{\"o}{\ss}en durchgef{\"u}hrt wird. Zudem bedeutet experimentell, dass Einflussfaktoren gezielt ver{\"a}ndert werden und beobachtet wird, wie sich dadurch eine abh{\"a}ngige Messgr{\"o}{\ss}e ver{\"a}ndert.

Das Versuchsdesign folgt somit der klassischen Struktur eines kontrollierten Experiments:

\begin{itemize}
\item \textbf{Unabh{\"a}ngige Variable (UV):} Die aktive Zephyr-Konfiguration (Minimal, Logging, IoT-Stack) ist der zu manipulierende Faktor.
\item \textbf{Abh{\"a}ngige Variablen (AV):} Die f{\"u}nf gemessenen RTOS-Leistungsmetriken (Context-Switch, Interrupt Latenz, Timer-Jitter, Heap-Allokationsgeschwindigkeit, Heap-Fragmentierung), wobei hierbei die Wirkung der Zephyr-Konfigurationen beobachtet wird.
\item \textbf{Kontrollierte Variablen:} Alle f{\"u}r die Messung relevanten Rahmenbedingungen, die zwischen den Durchl{\"a}ufen konstant gehalten werden, um Reproduzierbarkeit zu erm{\"o}glichen. Dazu geh{\"o}ren die verwendete Hardware, CPU-Taktfrequenz, Compiler, Zephyr-Version, Messverfahren, Stichprobenumfang und Umgebung. Die Anzahl der Messwiederholungen und der Messmechanismus f{\"u}r die Datenerhebung f{\"u}r alle Konfigurationen sind ebenfalls festgelegt.
\end{itemize}
Jede Kombination aus Konfiguration und Metrik (In der restlichen Arbeit als \glqq{}Zelle\grqq{} bezeichnet) wird als eigenst{\"a}ndiges, minimales Firmware-Image gebaut. Dabei w{\"a}hlt CMakeLists.txt durch die Methode \lstinline[style=fileinline]|target_sources_ifdef| genau eine \lstinline[style=fileinline]|bm_*.c|-Datei zur Build-Zeit aus. \lstinline[style=fileinline]|Kconfig| macht dar{\"u}ber hinaus die Auswahl {\"u}ber \lstinline[style=fileinline]|choice BENCH_TYPE| explizit. Die getrennten Firmware-Images verringern dabei die gegenseitige Beeinflussung durch gleichzeitig eingebundenen Benchmark-Code. Unterschiede im erzeugten Speicherlayout bleiben jedoch dennoch m{\"o}glich.

Die drei Konfigurationsstufen sind zudem kumulativ/ inkrementell angelegt. Das Logging-Profil entsteht, indem \lstinline[style=fileinline]|minimal.conf| und das Konfigurationsfragment \lstinline[style=fileinline]|logging.conf| gemeinsam eingebunden werden. Dabei enth{\"a}lt \lstinline[style=fileinline]|logging.conf| nur die Einstellungen f{\"u}r das Logging-Subsystem und bildet damit die Konfigurationsdifferenz im Vergleich zu der Minimal-Konfiguration. Dadurch wird nicht nur das IoT-Stack mit dem Kernel-Build verglichen, sondern auch der isolierte Effekt einzelner Subsysteme, in diesem Fall das Logging-Subsystem, dass sauber quantifizierbar bleibt. Durch diese Methodik wird eine Stufenanalyse erm{\"o}glicht anstelle eines Zwei-Gruppen-Vergleichs.

F{\"u}r jede Zelle des Versuchsplans wird eine Stichprobe von \(N\) Messwerten erhoben, aus der deskriptive Kennzahlen berechnet werden. Diese sind das Minimum, Maximum, der Mittelwert und die Stichproben-Standardabweichung. Au{\ss}erdem existiert ein Referenzpunkt gegen Zephyrs eigene Benchmarks, welche dieselbe Timing-API benutzen. Dies dient als externe Validierung der selbst entwickelten Messungen.

\section{Zielplattform}
Als Zielplattform f{\"u}r die Untersuchungen wird das nRF52840 Development Kit von Nordic Semiconductor verwendet. Die Plattform basiert auf einem 32-Bit Arm Cortex-M4 mit einer Taktfrequenz von 64 MHz und verf{\"u}gt {\"u}ber 1 MB Flash-Speicher sowie 256 kB RAM. Damit bietet der Mikrocontroller eine klar definierte und ressourcenbeschr{\"a}nkte Hardwarebasis f{\"u}r die Untersuchung des Laufzeit- und Speicherverhaltens von Zephyr RTOS.

F{\"u}r die Untersuchung der Laufzeiteigenschaften sind vor allem der Nested Vectored Interrupt Controller (NVIC) und der Data Watchpoint and Trace (DWT) des Cortex-M4 relevant. Der NVIC {\"u}bernimmt die Verarbeitung und Priorisierung von Interrupts und bildet damit die hardwareseitige Grundlage f{\"u}r die Messung der Interrupt-Latenz. Der DWT erm{\"o}glicht die Erfassung von CPU-Taktzyklen und wird f{\"u}r die hochaufl{\"o}sende Laufzeitmessung der Benchmarks eingesetzt. Bei 64 MHz entspricht ein CPU-Zyklus 15,625 ns, wodurch bspw. Kontextwechsel- und Interrupt-Laufzeiten mit hoher zeitlicher Aufl{\"o}sung bestimmt werden k{\"o}nnen.

F{\"u}r die Untersuchung des Timer-Jitters wird auf der anderen Seite die RTC-basierte Zeitbasis von Zephyr verwendet. In der verwendeten Konfiguration wird hierf{\"u}r RTC1 mit einem 32,768-kHz-Low-Frequency-Clock betrieben. Da diese Zeitbasis auch w{\"a}hrend des Prozessorschlafs weiterl{\"a}uft, eignet sie sich zur Messung periodischer Timerintervalle im Tickless-Betrieb. Die resultierende Aufl{\"o}sung betr{\"a}gt etwa 30,5 \textmu{}s pro Tick. Die Kombination aus Cortex-M4, Interrupt- und Zeitgeberhardware sowie den begrenzten Speicherressourcen bildet damit die technische Grundlage f{\"u}r die durchgef{\"u}hrten Benchmarks (Nordic Semiconductor, o. D., Abschnitt 6.20; Zephyr Project, 2024, nrf\_rtc\_timer.c, Tag v3.7.2).

\section{Softwareumgebung}
Das Build befindet sich auf demselben Ubuntu-Host im Arbeitsbereich \lstinline[style=fileinline]|~/zephyr-ws|. Dabei handelt es sich um eine eigene Python-venv (Virtual Environment). Das Flashen auf den Mikrocontroller erfolgt {\"u}ber denselben J-Link-OB-Debugprobe via \lstinline[style=fileinline]|west flash| (nrfjprog-Runner). Die Ausgabe erfolgt dann {\"u}ber denselben UART/ VCOM-Kanal. Pro Messung werden ebenfalls der Commit-Hash und das verwendete \lstinline[style=fileinline]|EXTRA_CONF_FILE| notiert.

\section{Unabh{\"a}ngige Variablen}
Die UV werden als Kconfig-Fragmente realisiert, das {\"u}ber \lstinline[style=fileinline]|EXTRA_CONF_FILE| zus{\"a}tzlich zur minimalen Basis \lstinline[style=fileinline]|prj.conf| der Mess-App eingebunden wird. Die Basis enth{\"a}lt nur das, was f{\"u}r jede Messung zwingend n{\"o}tig ist, darunter die Timing-API, UART/ printk-Ausgabe und das Heap f{\"u}r die Allocation-Benchmarks. Diese Basis ist bewusst identisch in allen drei Konfigurationsstufen, damit diese nicht Teil der Variation wird.

\subsection{Minimal Konfiguration}
Der schlanke RTOS-Kern ohne jegliche Zusatzdienste. Boot-Banner, Timeslicing, Assertions, Logging, Thread-Namen, MPU/ Stack-Guard, SEGGER-RTT und GPIO-Treiber sind explizit deaktiviert. Diese Konfigurationsstufe dient als Baseline und liefert die verringerten Kosten der untersuchten Kerneloperationen inklusive der im Messpfad verbleibenden Anteile. Besonders ohne MPU-bedingte Stack-Guard-Reprogrammierung bei jedem Context Switch und ohne durch Timeslicing verursachte Scheduler-Unterbrechungen.

\subsection{Logging}
Diese Konfiguration baut auf der Minimal-Konfiguration auf und aktiviert das Zephyr-Logging-Subsystem im Deferred-Modus. \lstinline[style=fileinline]|CONFIG_LOG_PRINTK=n| verhindert, dass die {\"u}ber \lstinline[style=fileinline]|printk()| ausgegebenen Benchmark-Ergebnisse in das Logging-Subsystem umgeleitet werden. BENCHMETA-, BENCH- und BENCHEND-Zeilen bleiben dadurch auf dem synchronen printk-/ UART-Pfad. Das Logging-Backend verwendet dabei ebenfalls UART. Die Benchmark-Implementierungen erzeugen jedoch selbst keine LOG\_*-Nachrichten. Gemessen werden daher vor allem Speicherbedarf, Initialisierung, vorhandene Logging-Strukturen und m{\"o}gliche indirekte Layout- oder Hintergrundeffekte, nicht die Kosten eines fortlaufenden Nachrichtenstroms.

\subsection{IoT-Stack}
Diese Konfiguration baut ebenfalls auf der Minimal-Konfiguration auf. Es wird zudem durch Netzwerk- und Sicherheits-Subsysteme erweitert, welche f{\"u}r einen IoT-Anwendungsfall typisch sind. Dabei handelt es sich um OpenThread als Minimal Thread Device (MTD), CoAP und eine DTLS-Absicherung {\"u}ber einen Pre-Shared Key (siehe 2.7). MCUmgr und MCUboot sind hierbei bewusst nicht Teil der Konfigurationsstufe, da diese am Sysbuild-Mechanismus h{\"a}ngen und daher nichts zur Frage beitragen, wie sich Netzwerk- und Sicherheits-Subsysteme die f{\"u}nf Laufzeitmetriken beeinflussen. Sie werden daher in der Anwendung \lstinline[style=fileinline]|iot_sensor| betrachtet (siehe 6.4). Zudem wurde {\"u}ber \lstinline[style=fileinline]|CONFIG_OPENTHREAD_MANUAL_START=y| der Thread konfiguriert, damit die Benchmarks nicht durch endlose Netzwerk-Beitrittsversuche ohne vorhandenen Thread-Koordinator gest{\"o}rt werden. Durch \lstinline[style=fileinline]|CONFIG_OPENTHREAD_MANUAL_START=y| wird kein automatischer Netzwerkbeitritt gestartet. Der OpenThread-Netzwerkbetrieb bleibt w{\"a}hrend der Benchmarks inaktiv. Erfasst werden dabei das aktivierte Softwareprofil, sein statischer Speicherbedarf, Initialisierung und vorhandene Kernelobjekte. Der aktive Mesh-Verkehr und ein DTLS-Handshake werden hierbei nicht gemessen.

\section{Abh{\"a}ngige Variablen}
Die DWT-basierten Benchmarks f{\"u}r Kontextwechsel, Interrupt-Latenz und Heap-Allokation benutzen Zeitstempel {\"u}ber \lstinline[style=fileinline]|bench_now()|, einen kalibrierten Abzug des Messaufwands und \lstinline[style=fileinline]|struct bench_stats|. Der Timer-Benchmark benutzt dieselbe Statistik und Ergebnisausgabe, liest jedoch RTC-Systemtimerzyklen {\"u}ber \lstinline[style=fileinline]|sys_clock_cycle_get_32()|. Die Heap-Fragmentierung ist keine Zeitmessung und besitzt deshalb keinen Zeitstempel und keinen Abzug des Messaufwands oder Welford-Statistik.

\subsection{Context-Switch-Latenz}
Der Context-Switch wurde im Benchmark \lstinline[style=fileinline]|bm_ctx_switch.c| implementiert. Hier wird die Zeit vom Ausl{\"o}sen eines Threadwechsels bis zum Zeitstempel als erster C-seitiger Aktion im Zielthread gemessen. Dies erfolgt in zwei Varianten:

\lstinline[style=fileinline]|ctx_switch_sem_wakeup|: Semaphore-geweckter h{\"o}herpriorer Worker-Thread

\lstinline[style=fileinline]|ctx_switch_yield|: Kooperatives Ping-Pong von gleichprioren Threads via \lstinline[style=fileinline]|k_yield()|.

Dieser Benchmark erfasst dabei die Scheduler-Entscheidung inklusive PendSV-Kontextwechsel, die Grundkosten jeder Multithreading-Architektur. Dies bildet die zentrale Vergleichsgr{\"o}{\ss}e zwischen den drei Konfigurationen.

\subsection{Interrupt-Latenz}
Die Interrupt-Latenz wurde im Benchmark \lstinline[style=fileinline]|bm_irq_latency.c| implementiert. Es wird die Zeit vom Software-Trigger bis zum Zeitstempel am Anfang des registrierten C-Handlers gemessen. Der Exception-Eintritt, Zephyr-Wrapper und Compiler-Prolog liegen schon davor und sind Teil des Messwerts. Dies erfolgt durch \lstinline[style=fileinline]|NVIC_SetPendingIRQ| auf einer ungenutzten Peripherie-Leitung des nrf52840. Dieser Benchmark zeigt, wie stark Kernel-Funktionen, wie bspw. Logging, den Determinismus-kritischsten Pfad eines RTOS beeinflussen.

\subsection{Timer-Jitter}
Der Benchmark \lstinline[style=fileinline]|bm_timer_jitter.c| erfasst die Abst{\"a}nde zwischen aufeinanderfolgenden Z{\"a}hlerzugriffen im Callback eines periodischen 10-ms-\lstinline[style=fileinline]|k_timer|. Daf{\"u}r wird \lstinline[style=fileinline]|sys_clock_cycle_get_32()| verwendet. Die Ausgabe hei{\ss}t \lstinline[style=fileinline]|timer_period_rtc_cycles| und besitzt die Einheit RTC-Systemtimerzyklen. \lstinline[style=fileinline]|TIMERMETA| gibt die Frequenz und die mit \lstinline[style=fileinline]|k_ms_to_cyc_ceil32()| berechnete Periode an. Der Benchmark gibt Periodenwerte aus. Die Abweichung und Streuung im Vergleich zur Sollperiode werden daraus ausgewertet.

\subsection{Heap-Allokationsgeschwindigkeit}
Die Heap-Allokationsgeschwindigkeit wurde im Benchmark \lstinline[style=fileinline]|bm_heap_alloc.c| implementiert. Hier wird die Zeit f{\"u}r die Funktionen \lstinline[style=fileinline]|k_malloc| und \lstinline[style=fileinline]|k_free| eines 64-Byte-Blocks gemessen. Allokation und Freigabe folgen direkt aufeinander. Dadurch bleibt der Systemheap nahezu leer und unfragmentiert, sodass ein vorteilhafter Ausgangszustand untersucht wird. Dieser Benchmark ist relevant, weil \lstinline[style=fileinline]|sys_heap| eine variable Laufzeit hat, was f{\"u}r RTS kritisch ist.

\subsection{Heap-Fragmentierung}
Die Heap-Fragmentierung wurde im Benchmark \lstinline[style=fileinline]|bm_heap_frag.c| implementiert. Hierbei handelt es sich nicht um ein Zeit-, sondern um ein Zustandsma{\ss}. Ein eigener 8-KiB k\_heap wird mit 32x128-Byte-Bl{\"o}cken gef{\"u}llt. Dabei wird im Schachbrettmuster jeder zweite freigegeben und anschlie{\ss}end die angeforderten Nutzbytes, die rechnerisch nicht angeforderten Bytes und die gr{\"o}{\ss}te im 8-Byte-Suchraster erfolgreich gepr{\"u}fte Einzelanforderung protokolliert. Dieser Freiblock wird durch eine eigene Bisektions-Probe protokolliert, da \lstinline[style=fileinline]|sys_heap| diese Kennzahl nicht direkt liefert. Dieser Benchmark zeigt ein f{\"u}r MMU-lose Systeme spezifisches Risiko: Rechnerisch ausreichender, aber durch Fragmentierung nicht mehr am St{\"u}ck nutzbarer Speicher.

Kontextwechsel, Interrupt-Latenz und Heap-Allokation liefern CPU-Zyklen. Der Timer-Benchmark liefert RTC-Systemtimerzyklen. Die jeweilige Einheit ergibt sich dabei aus dem Metriknamen und \lstinline[style=fileinline]|TIMERMETA|. Die Heap-Fragmentierung liefert Nutzbyte-Buchf{\"u}hrungswerte und Probe-Ergebnisse in Byte und benutzt deshalb \lstinline[style=fileinline]|HEAPSTATS| und \lstinline[style=fileinline]|HEAPPROBE|.

\section{Kontrollierte Variablen}
Diese Variablen werden {\"u}ber alle Konfigurationsstufen und Metriken hinweg konstant gehalten, um zu verhindern, dass Effekte f{\"a}lschlich der UV zugeschreiben werden.

\subsection{Hardware/ CPU-Frequenz}
Ein einziges physisches Board, das Nordic nrf52840 DK, mit einer CPU basierend auf der Cortex-M4F Architektur und einer Takt-Frequenz von 64 MHz. Es verf{\"u}gt {\"u}ber 1 MB Flash und 256 KB RAM. Die Frequenz wird dar{\"u}ber hinaus zur Laufzeitkontrolle bei jedem Boot {\"u}ber die Methode \lstinline[style=fileinline]|timing_freq_get()| in der \lstinline[style=fileinline]|BENCHMETA|-Zeile in \lstinline[style=fileinline]|main.c| mit ausgegeben. Falls es also zu einer Abweichung in der Takt-Frequenz kommen sollte, w{\"a}re dies sofort sichtbar.

\subsection{Compiler und Werkzeugkette}
In dieser Untersuchung wurde das Zephyr Software Development Kit (SDK) 0.16.9 verwendet. Der Compiler ist daher GCC 12.2.0, mit Binutils 2.38 und der Picolibc 1.8.6. Diese wurden {\"u}ber ZEPHYR\_SDK\_INSTALL\_DIR gepinnt, sodass sie identisch f{\"u}r alle Builds sind.

\subsection{Compiler optimierung}
Keines der drei Konfigurationsfragmente {\"a}ndert das Optimierungsniveau. Vor der Ergebnisinterpretation muss durch Compileraufrufe best{\"a}tigt und dokumentiert werden, dass alle Builds mit derselben Zephyr-Optimierung erstellt wurden.

\subsection{Zephyr Version}
F{\"u}r diese Untersuchung wird Zephyr v3.7.2 aus der Long-Term-Support-Reihe 3.7 verwendet. Dies wurde im Manifest \lstinline[style=fileinline]|west.yml| {\"u}ber revision: 3.7.2 mit \lstinline[style=fileinline]|import: true| gepinnt. Dies {\"u}bernimmt passende HAL-, MCUBoot-, mbedTLS- und OpenThread-Revisionen. Das gepinnte Manifest legt hierbei den vorgesehenen Versionsstand fest. F{\"u}r die Reproduzierbarkeit ist dar{\"u}ber hinaus zu dokumentieren, welcher Stand wirklich ausgecheckt und gebaut wurde.

\subsection{Build-Konfiguration (Sockel)}
Die in \lstinline[style=fileinline]|prj.conf| definierte Basis ist f{\"u}r alle drei UV-Stufen identisch. Diese beinhaltet \lstinline[style=fileinline]|Benchlib, printk/ UART/ Konsole, CONFIG_CBPRINTF_FP_SUPPORT| f{\"u}r Flie{\ss}komma-Ausgabe der Statistik und einen 64-KiB-System-Heap. Diese Basis wird nicht durch die Konfigurationen ver{\"a}ndert, sondern lediglich erg{\"a}nzt. Dadurch bleibt bspw. der ROM-Messaufwand der Flie{\ss}komma-Ausgabe in jeder Konfiguration gleich gro{\ss} und erm{\"o}glicht somit einen Vergleich.

\subsection{Messmechanismus (Timing/ Statistik)}
Die kurzen Laufzeitmessungen benutzen Zephyrs Timing-API und auf dem Cortex-M4F den DWT-Zyklenz{\"a}hler. Zwei direkt aufeinanderfolgende Reads werden 100-mal gemessen. Das kleinste beobachtete Delta wird von den DWT-basierten Ergebnissen abgezogen. Der Timer-Benchmark benutzt stattdessen den RTC1-basierten Systemtimer und keinen Abzug vom DWT-Messaufwand.

\[t_i = \frac{c_i}{64\,000\,000}\,\mathrm{s}\]
\[t_i = \frac{c_i}{32\,768}\,\mathrm{s}\]
\[t_i = \frac{c_i}{32\,768}\,\mathrm{s}\]
Die Welford-Statistik berechnet Minimum, Maximum, Mittelwert und Stichproben-Standardabweichung f{\"u}r die wiederholten Zeit- bzw. Periodenwerte. Heap-Fragmentierung und statischer Speicherbedarf werden dabei nicht mit dieser Statistik ausgewertet.

\subsection{Anzahl der Messungen}
F{\"u}r die wiederholten DWT- und Timer-Messreihen werden 1000 Werte nach zehn verworfenen Warm-up-Perioden bzw. -Durchl{\"a}ufen ausgewertet. Diese Parameter werden {\"u}ber \lstinline[style=fileinline]|CONFIG_BENCHLIB_ITERATIONS| und \lstinline[style=fileinline]|CONFIG_BENCHLIB_WARMUP| dokumentiert. Heap-Fragmentierung und RAM-/ ROM-Bedarf sind auf der anderen Seite phasenbezogene Zustands- bzw. Build-Messungen und besitzen daher keine Stichprobe aus 1000 Wiederholungen.

\section{Ablauf einer Messung}
Dieser Abschnitt beschreibt den kompletten Ablauf einer Messung, welcher f{\"u}r jede Messzelle gleich ist. Dabei wird kurz gezeigt, wie der Arbeitsbereich eingerichtet wird. Die Werkzeugkette wird ebenfalls einmal eingerichtet. Das eigentliche Bauen, Flashen, Auslesen und die Dokumentationspflicht sind dabei bewusst als ein fester und wiederholbarer Ablauf gestaltet, damit jede Messung genau nachvollzogen werden kann.

\textbf{Einmalige Einrichtung }Vor der ersten Messung wurde der West-Arbeitsbereich einmalig durch das offizielle Zephyr-Manifest in der Version v3.7.2 initialisiert. Dies wurde dabei durch das eigene Projektmanifest \lstinline[style=fileinline]|thesis/west.yml| erweitert (siehe 6.1).

\begin{lstlisting}[style=basecode,language=BAsh]
mkdir -p ~/zephyr-ws && cd ~/zephyr-ws
python3 -m venv .venv
source .venv/bin/activate
pip install west
west init --mr v3.7.2
west update
pip install -r zephyr/scripts/requirements.txt
west zephyr-export
west config manifest.path thesis
west update
\end{lstlisting}
Anschlie{\ss}end wurden das Zephyr SDK 0.16.9 und die f{\"u}r das Flashen gebrauchten Werkzeuge, darunter J-Link und nRF Command Line Tools, installiert. Diese wurden zudem {\"u}ber einen Smoke-Test mit \lstinline[style=fileinline]|zephyr/samples/basic/blinky| verifiziert.

\textbf{Beginn jeder Sitzung }Am Anfang jeder Mess-Sitzung wird in den Arbeitsbereich gewechselt, die Python-virtuelle Umgebung aktiviert und der aktuell ausgecheckte Commit-Hash des eigenen Projekts notiert.

\begin{lstlisting}[style=basecode,language=BAsh]
cd ~/zephyr-ws
source .venv/bin/activate
git -C thesis pull
git -C thesis log --oneline -1
\end{lstlisting}
Dieser Commit-Hash wird dabei jeder in dieser Arbeit dokumentierten Messung zugeordnet, damit sich jede Messreihe genau auf den zum Messzeitpunkt verwendeten Software- und Konfigurationsstand zur{\"u}ckf{\"u}hren l{\"a}sst (siehe Kapitel 4.3).

\textbf{Baseline pro Konfigurationsstufe }F{\"u}r jede der drei Konfigurationsstufen wird zuerst, noch vor den eigentlichen f{\"u}nf Benchmarks, ein \lstinline[style=fileinline]|CONFIG_BENCH_NONE-Build| erzeugt. Dieser Build enth{\"a}lt keine \lstinline[style=fileinline]|bm_*.c|-Datei, aber weiterhin \lstinline[style=fileinline]|main.c|, \lstinline[style=fileinline]|benchlib|, Timing-Initialisierung, UART-/ printk-Ausgabe, Flie{\ss}kommaformatierung und den konfigurierten 64-KiB-Systemheap. Er dient als gemeinsamer Mess-Harness zur Bestimmung des statischen Speicherbedarfs des jeweiligen Profils ohne benchmarkspezifischen Testcode. Da \lstinline[style=fileinline]|rom_report| und \lstinline[style=fileinline]|ram_report| nur die Linker-Map des gebauten ELF-Abbilds auswerten, ist f{\"u}r diesen Schritt kein Flashen und keine angeschlossene Hardware n{\"o}tig.

\begin{lstlisting}[style=basecode,language=BAsh]
cd ~/zephyr-ws
west build -p always -b nrf52840dk/nrf52840 thesis/apps/bench -- \
  -DEXTRA_CONF_FILE=../../configs/minimal.conf -DCONFIG_BENCH_NONE=y
west build -t rom_report | tail -40
west build -t ram_report | tail -40
\end{lstlisting}
F{\"u}r die Logging-Konfiguration wird dabei \lstinline[style=fileinline]|minimal.conf| zusammen mit dem Fragment \lstinline[style=fileinline]|logging.conf| {\"u}bergeben. F{\"u}r die IoT-Stack-Konfiguration wird dazu noch \lstinline[style=fileinline]|iot_stack.conf| und das Devicetree-Overlay \lstinline[style=fileinline]|iot_stack.overlay| f{\"u}r die OpenThread-Entropiequelle {\"u}bergeben. Genau derselbe Baseline-Prozess wird also f{\"u}r alle drei Konfigurationsstufen gleich wiederholt. Nur das {\"u}bergebene \lstinline[style=fileinline]|EXTRA_CONF_FILE| unterscheidet sich.

\textbf{Ausf{\"u}hrung der f{\"u}nf Benchmarks }Nach der Baseline wird f{\"u}r dieselbe Konfigurationsstufe nacheinander jeder der f{\"u}nf Benchmarks gebaut, geflasht und ausgelesen. Dies geschieht dabei jeweils durch das Setzen eines \lstinline[style=fileinline]|CONFIG_BENCH_*|-Symbols aus der \lstinline[style=fileinline]|choice|-Auswahl (siehe 4.1).

\begin{lstlisting}[style=basecode,language=BAsh]
west build -p always -b nrf52840dk/nrf52840 thesis/apps/bench -- \
  -DEXTRA_CONF_FILE=../../configs/minimal.conf -DCONFIG_BENCH_CTX_SWITCH=y
west flash
\end{lstlisting}
Dieser Doppelschritt aus Build und \lstinline[style=fileinline]|west flash| wird zudem f{\"u}r \lstinline[style=fileinline]|CONFIG_BENCH_IRQ_LATENCY|, \lstinline[style=fileinline]|CONFIG_BENCH_TIMER_JITTER|, \lstinline[style=fileinline]|CONFIG_BENCH_HEAP_ALLOC| und \lstinline[style=fileinline]|CONFIG_BENCH_HEAP_FRAG| wiederholt. Dazu wird vor der eigentlichen Messreihe einer Konfigurationsstufe einmalig CONFIG\_BENCH\_SELFTEST gebaut und geflasht, um die Kalibrierungspr{\"u}fung der Zeitmesskette zu verifizieren (siehe Kapitel 5). F{\"u}r die Logging- und IoT-Stack-Konfiguration wird derselbe Prozess von sechs Bauvorg{\"a}ngen (Selbsttest und die f{\"u}nf Benchmarks) mit dem jeweils passenden \lstinline[style=fileinline]|EXTRA_CONF_FILE|- bzw. extra \lstinline[style=fileinline]|EXTRA_DTC_OVERLAY_FILE|-Parameter wiederholt. Insgesamt besteht dies f{\"u}r jede der drei Konfigurationsstufen aus sieben Bauvorg{\"a}ngen: Die Baseline (\lstinline[style=fileinline]|BENCH_NONE|), den Selbsttest und die f{\"u}nf Benchmarks.

\textbf{Auslesen }Nach jedem Flashvorgang wird die Ausgabe {\"u}ber die serielle Schnittstelle des J-Link-Onboard-Debuggers ausgelesen.

\begin{lstlisting}[style=basecode,language=BAsh]
picocom -b 115200 /dev/ttyACM0
\end{lstlisting}
Nach dem Dr{\"u}cken der Reset-Taste am Board erscheint die Ausgabe der jeweils aktiven Firmware. Vor jeder Auswertung wird dabei kontrolliert, ob das Build-Log f{\"u}r jedes benutzte Konfigurationsfragment eine Zeile in der Form \lstinline[style=fileinline]|Merged configuration '.../<name>.conf'| und \lstinline[style=fileinline]|Found toolchain: zephyr 0.16.9| enth{\"a}lt. Au{\ss}erdem wird gepr{\"u}ft, ob bei der Minimal- und Logging-Konfiguration kein Boot-Banner beim Reset zu sehen ist. Ein Banner w{\"u}rde anzeigen, dass \lstinline[style=fileinline]|minimal.conf| nicht richtig angewendet wurde. Eine Beispiel-Ausgabe des Selbsttests sieht wie folgt aus:

\begin{lstlisting}[style=basecode,language={}]
BENCHMETA,zephyr=3.7.2,freq_hz=64000000,overhead_cycles=12,iterations=1000,warmup=10
BENCH,selftest_busy_wait_100us,1000,6423,6435,6434.96,0.57
BENCHEND
\end{lstlisting}
\textbf{Dokumentation }F{\"u}r jede Messzelle des Versuchsplans werden zuletzt das Datum, der Commit-Hash, die benutzte Konfiguration inklusive des/ der \lstinline[style=fileinline]|EXTRA_CONF_FILE|-Fragments/-e, der aktive Benchmark, die komplette BENCHMETA-Kopfzeile, die dazugeh{\"o}rigen \lstinline[style=fileinline]|BENCH|-, \lstinline[style=fileinline]|HEAPSTATS|- bzw. \lstinline[style=fileinline]|HEAPPROBE|-Zeilen notiert. Au{\ss}erdem wird bei Fu{\ss}abdruck-Messungen die \glqq{}Total\grqq{}-Zeile von \lstinline[style=fileinline]|rom_report| und \lstinline[style=fileinline]|ram_report| dokumentiert. Dadurch kann bei jeder Kennzahl (siehe Kapitel 7) bis zur Rohausgabe der Hardware zur{\"u}ckverfolgt werden.
